\scanpage{96}

three, compatible with one and the same homomorphism
\diagram{\begin{tikzcd}[ampersand replacement=\&,column sep=large]
\mathcal M_{\deleted{\unsure{0}}} \arrow[r] \&
\deleted{\unsure{\mathbb O|\Phi}}\gamma_{\rho,\sigma}
\overset{\unsure{\mathrm{def}}}{=}
\underbrace{\unsure{\underline{\gamma}_{\rho,\sigma}^{*}|\underline{\gamma}_{\rho,\sigma}^{\bullet}}}_{\mathbb O_{12}^{*}}
(\widehat{\mathbb Z})^{\Delta}
=\deleted{\unsure{\mathbb O_{12}^{*}}}\widehat{\mathbb Z}^{*}
\end{tikzcd}}
which is none other than the cyclotomic character $\chi$.
\unsure{One}\deleted{\unsure{the|of}} of these three,
\diagram{\begin{tikzcd}[ampersand replacement=\&,column sep=small,row sep=small]
\mathcal M \arrow[rr] \arrow[dr,dashed] \&\& N_{\rho}^{\sim}\arrow[r,"\sim"] \&\widehat{\mathbb Z}^{*},\\
\&\inserted{\Gamma_{\mathbb Q}^{\mathrm{et}}}\arrow[ur,dashed]\&\&
\end{tikzcd}}
reduces here \unsure{to the said} character, and thus raises no interesting question. We have on the other hand

\noindent (218)\quad
\diagram{\begin{tikzcd}[ampersand replacement=\&,column sep=large]
\mathcal M\arrow[r]\arrow[d]\&G_{\rho,\sigma}=\widehat{\mathbb Z}^{*}\mathbin{\cdot}\mathfrak S^{+\wedge}\arrow[d]\\
\mathcal M^{!}\arrow[r]\&G_{\rho,\sigma}^{\unsure{\natural|\sharp}}=\widehat{\mathbb Z}^{*}\mathbin{\cdot}M^{\wedge},
\end{tikzcd}}

which will have to be examined more closely --- as the cyclotomic character is \inserted{(even on $\mathcal M!$)} surjective, we \unsure{know} in any case that these two homomorphisms are \uline{surjective}. Finally, we have a homomorphism
$\mathcal M\longrightarrow N_{\sigma}^{\unsure{\natural|\sharp}\deleted{\unsure{\mathrm{ab}}}}$, which factors as

\noindent (219)\quad
\diagram{\begin{tikzcd}[ampersand replacement=\&,column sep=large]
\Gamma_{\mathbb Q}^{\mathrm{et}}\arrow[r]\&N_{\sigma}^{\unsure{\natural|\sharp}\deleted{\unsure{\mathrm{ab}}}}=\widehat{\mathbb Z}^{*}\mathbin{\cdot}\widehat{\mathbb Z}
\end{tikzcd}}
\hfill ($\frac12$ direct)

which will also have to be examined --- in particular see whether it is surjective.
Of course, I have a chance of doing so only by the ``arithmetic route'', by \unsure{proving} that the homomorphism induced on $\Gamma_{\mathbb Q}$ itself is an epimorphism.

\margin{95}

\scanpage{97}
\mspage{696}

I feel like putting these homomorphisms in relation with the ``tower'' of algebraic curves defined over $\mathbb Q$, indexed by $n\in\unsure{\mathbb N^{*}|\mathbb N^{\wedge}}$

\noindent (220)\quad
\diagram{\begin{tikzcd}[ampersand replacement=\&,column sep=small,row sep=small]
X_n : \arrow[d,phantom,"\subset" description,sloped] \& x_0^n+x_1^n+x_{\infty}^n=0\\
\mathbb P_{\mathbb Q}^{2}\&
\end{tikzcd}}
\hfill (Fermat curve of index $n$)

\inserted{Let $U_n$ be the open set of points where $x_0x_1x_{\infty}\ne0$}

On $X_n$ the group scheme $\mu_n^{\{0,1,\infty\}}/\underbrace{\mu_n}_{\mathrm{diag}.}$ operates ($\mu_n={}_n\mathbb G_m$), \deleted{and \unsure{we have}} as does $\mathfrak S_3=\mathfrak S_{\{0,1,\infty\}}$, hence also the semidirect product
\[
(221)\qquad \underline F_n=\mathfrak S_3\mathbin{\cdot}\underbrace{\mu_n^{\{0,1,\infty\}}/\mu_n}_{\underline F_n^{\circ}},
\]
\margin{\unsure{Notation} $F_n$ not very \unsure{happy} --- it is possible that there will be need later for the letter $F_n$, $F_{\infty}$ for the Fermat coverings $U_n$, $U_{\infty}$ themselves\unsure{\ldots} (compare \unsure{corrections|rectifications} page~719!)}
we have
\[
(222)\qquad X_n/\underline F_n^{\circ}\simeq X_1\simeq U_{0,3},
\]
the isomorphism compatible with the actions of $\mathfrak S_3$. More generally, we have, if $n=n'd$, a morphism

\noindent (223)\quad
\diagram{\begin{tikzcd}[ampersand replacement=\&,column sep=large,row sep=small]
X_{\unsure{n=n'd}}\arrow[r,very thick]\&X_{n'}\\
x_i\arrow[r,mapsto]\&x_i^d
\end{tikzcd}}

and in this way $\deleted{X_n}\inserted{U_n}$ becomes a principal covering (= torsor) over $\deleted{X_{n'}}\inserted{U_{n'}}$, with group $\mu_d^{\{0,1,\infty\}}/\mu_d$, More precisely still, \deleted{(223) is compatible with}

\noindent \deleted{(224)}\quad
\diagram{\begin{tikzcd}[ampersand replacement=\&,column sep=large,row sep=small]
\deleted{F_n}\arrow[r,"\deleted{\text{epimorphism}}","\deleted{\unsure{\text{canonical}}}"']\arrow[d,equal]\&\deleted{F_{n'}}\arrow[d,equal]\\
\deleted{\mathfrak S_3\mathbin{\cdot}\mu_n^{\{0,1,\infty\}}}\&\deleted{\mathfrak S_3\mathbin{\cdot}\mu_{n'}^{\{0,1,\infty\}}}
\end{tikzcd}}
\deleted{$\left\{\begin{aligned}g&\longmapsto g\quad\text{if }g\in\mathfrak S_3,\\(x_0,x_1,x_{\infty})&\longmapsto (x_0^d,x_1^d,\unsure{x_{\infty}^d})\end{aligned}\right.$}

\margin{96}

\scanpage{98}

\deleted{and $\mu_n^{\{0,1,\infty\}}$, a fortiori $\mu_d^{\{0,1,\infty\}}$, operates freely on $X_n$, and we find by passage to the quotient} \unsure{we have} a canonical exact sequence

\noindent (224)\quad
\diagram{\begin{tikzcd}[ampersand replacement=\&,column sep=small,row sep=small]
1\arrow[r,no head]\&\mu_d^{\{0,1,\infty\}}\deleted{/\mu_d}\arrow[r]\&
\begin{gathered}\underline F_n\\ \Vert\\ \unsure{\mathfrak S_{3\,S_0}}\mathbin{\cdot}\mu_n^{\{0,1,\infty\}}\end{gathered}\arrow[r]\&
\begin{gathered}\underline F_{n'}\\ \Vert\\ \mathfrak S_{3\,S_0}\mathbin{\cdot}\mu_{n'}^{\{0,1,\infty\}}\end{gathered}\arrow[r]\&1\\
\&(x_0,x_1,x_{\infty})\arrow[r,mapsto]\&(x_0,x_1,x_{\infty})\&\&\\
\&\&(x_0,x_1,x_{\infty})\arrow[r,mapsto]\&(x_0^d,x_1^d,x_{\infty}^d)\&\\
\&\&g\deleted{\unsure{\in\mathfrak S_3}}\arrow[r,mapsto]\&g\&
\end{tikzcd}}
\[
\text{for}\quad\begin{aligned}x_0,x_1,x_{\infty}&\in\Gamma\mu_n,\\g&\in\Gamma\mathfrak S_{3\,S_0}\end{aligned}\qquad\big)
\]

and \unsure{the morphism|the isomorphism} (223) is compatible with (224), and induces an isomorphism (compatible with actions of $F_{n'}$)
\[
(225)\qquad X_{n'}\simeq X_n/(\mu_d^{\{0,1,\infty\}}/\mu_d),
\]
Moreover, the action of $\deleted{\unsure{\mu_n^{\{0,}}}$
\[
\underline F_n^{\circ}=\mu_n^{\{0,1,\infty\}}/\mu_n
\]
on $\inserted{U_n}\deleted{X_n}$ is free, and a fortiori that of $\mu_d^{\{0,1,\infty\}}/\mu_d$ is, so that $\deleted{X_n}\inserted{U_n}$ is indeed a torsor over $\deleted{X_{n'}}\inserted{U_{n'}}$ with group
\[
\underline F_d^{\circ}=\mu_d^{\{0,1,\infty\}}/\mu_d
\]
Passing to the limit, we find
\[
(226)\qquad\left\{\begin{aligned}
X_{\infty}&=\varprojlim_{n\in\mathbb N^{*}}X_n,\quad U_{\infty}=\varprojlim U_n\\
\underline F_{\infty}&=\varprojlim\underline F_n=(\mathfrak S_3)_{S_0}\mathbin{\cdot}\deleted{\unsure{\widehat{\mathbb Z}^{2}}}\left(\widehat{\mathbb Z}^{\{0,1,\infty\}}/\widehat{\mathbb Z}\right)\otimes_{\widehat{\mathbb Z}}\unsure{T_{\bullet}(\mathbb G_m)|T(\mathbb G_m)}
\end{aligned}\right.
\]

\margin{97}

\scanpage{99}
\mspage{697}

with operation of $F_\infty$ on \deleted{$X_\infty$}\inserted{$U_\infty$}
\deleted{$F_\infty^\circ$ operates}
\[
(227)\qquad F_\infty^\circ=\varprojlim_n F_n^\circ\simeq
\inserted{\widehat{\mathbb Z}^{\{0,1,\infty\}}\otimes_{\widehat{\mathbb Z}}}
\unsure{T_\bullet|T_*}(\mathbb G_m)^{\deleted{\unsure{\{0,1,\infty\}}}}
\]
operating freely \inserted{on $U_\infty$} --- the quotient \unsure{being}
\deleted{$X_1$}\inserted{$U_1$}$\simeq U_{0,3}$.

\deleted{By the Galois theory of $(U_{0,3},\mathfrak S_3)$,}
\inserted{The pro-étale group $\underline F_\infty$ over $\mathbb Q$}
\deleted{this covering defines} corresponds \deleted{to a}
\inserted{profinite group on which $\Gamma_{\mathbb Q}$ operates, it is}
\deleted{of a homomorphism \unsure{defined} \illegible{2} \unsure{independent}}
\diagram{\begin{tikzcd}[ampersand replacement=\&,column sep=large]
\deleted{(227)} \& \deleted{\mathcal M_{0,3}=\Pi_1(U_{0,3\,\mathbb Q},\mathfrak S_3;-\bar j)} \arrow[r] \& {}
\end{tikzcd}}
the group
\[
(228)\qquad F_\infty^\circ\overset{\unsure{\mathrm{def}}}{=}
\underline F_\infty^\circ(\overline{\mathbb Q})\simeq
\deleted{\unsure{\mathfrak S_3}}\,\cdot\,
\overbrace{\left(
\underbrace{\widehat{\mathbb Z}^{\deleted{\unsure{*}}\{0,1,\infty\}}/\widehat{\mathbb Z}}_{\deleted{\unsure{F_\infty^\circ}}}
\right)}^{\unsure{\widehat M|M^*}}
\otimes_{\widehat{\mathbb Z}}
\underbrace{\widehat{\mathbb Z}}_{T_\bullet(\mathbb G_m)(\mathbb C)}
\]
on which $\Gamma_{\mathbb Q}$ operates \deleted{\unsure{via the trivial action}}
\deleted{on $\mathfrak S_3$, and}
\inserted{via} (the \deleted{\unsure{anti.}} cyclotomic character
\inserted{operating as a factor}
\deleted{on $\widehat{\mathbb Z}^{\unsure{\{0,1,\infty\}}}$}.
We \deleted{\unsure{the profinite group} \illegible{1}}
\inserted{make $\mathfrak S_3$ operate on the inverse
image of this group over $U_{0,3\,\mathbb Q}$,}
\deleted{by the action}
\inserted{in a way compatible with the action of $\mathfrak S_3$ on $U_{0,3\,\mathbb Q}$,}
\deleted{\illegible{2} over $\mathbb Q$ and \unsure{$\widehat{\mathbb Z}$} of \unsure{$U_{0,3}$} \illegible{1}}
\inserted{by permutation of the factors, and --- this corresponds
to an operation of \deleted{\unsure{$\mathcal M_{0,3}$}} $\mathcal M_{0,3}=\Pi_1(U_{0,3\,\mathbb Q},\mathfrak S_3;-\bar j)$}
\deleted{\unsure{permits} \unsure{to the fiber} \illegible{4}}
\inserted{on $F_\infty^\circ$ --- which is done via the canonical homomorphisms}
\diagram{\begin{tikzcd}[ampersand replacement=\&,column sep=large]
(229) \& \mathcal M \arrow[r] \& \mathfrak S_3, \& \mathcal M \arrow[r,"\chi"] \& \widehat{\mathbb Z}^{*}.
\end{tikzcd}}
\deleted{\unsure{$X_\infty$} \illegible{2}}

\margin{\uline{N.B.} The action of $\unsure{\mathcal M_{0,3}}$ on $\widehat M\simeq\underline M(\widehat{\mathbb Z})$ is none other than that induced by
\diagram{\begin{tikzcd}[ampersand replacement=\&,column sep=small]\mathcal M \arrow[r] \& \unsure{G}_{\rho,\sigma}\end{tikzcd}},
since $\mathcal M$ operates on the invariant subgroup $\widehat M$ of $G_{\rho,\sigma}$ --- \deleted{this} the action of $\omega$ being trivial, \unsure{is} factored into an action of \deleted{$\unsure{\mathcal M_{0,3}^{\mathrm{et}}}$}\inserted{$\mathcal M_{0,3}^{\unsure{\mathrm{et}}}$}.}

\[
(230)\qquad
\deleted{X_\infty}\inserted{U_\infty}(-\bar j)
\]
a torsor under $F_\infty^\circ$, on which
$\mathcal M_{0,3}=\Pi_1(U_{0,3\,\mathbb Q},\mathfrak S_3;-\bar j)$ operates
in a way compatible with its operations on $F_\infty^\circ$.

The datum of this profinite $(F_\infty^\circ,\mathcal M_{0,3})$-torsor
is equivalent to that of the principal covering
\deleted{$X_\infty$}\inserted{$U_\infty$} of $\unsure{\widehat U}_{0,3\,\mathbb Q}$,
\deleted{with group} $\underline F_\infty^\circ$,

\margin{98}

\scanpage{100}

with the operations of $\mathfrak S_3$ on it compatible
with the operations of $\mathfrak S_3$ on $\underline F_\infty^\circ$.

\deleted{On the other hand, the category of torsors under abelian groups
\deleted{$\unsure{\widehat{\mathbb Z}^{\{0,1,\infty\}}}$}
of the form $M\otimes_{\mathbb Z}T$, where $T$ is an invertible $\widehat{\mathbb Z}$-module,
is equivalent to that of torsors under the \unsure{integral} affine group
\[
\widehat{\mathbb Z}^{*}\cdot\widehat M
\]
}

\deleted{The datum of an isomorphism class of such torsors, with operation of
\deleted{$\unsure{\Gamma}$} $\mathcal M_{0,3}$ on them compatible with the action
on $\widehat M$ \unsure{(which|(by}}

Let us consider on the other hand the exact sequence
\diagram{\begin{tikzcd}[ampersand replacement=\&,column sep=small]
(231) \& 1 \arrow[r,dash] \& \underline M \arrow[r] \&
\underline G_{\rho,\sigma}/\unsure{\underline\mu_G|\mu_G}
\arrow[r,"{(\underline\chi,\mathrm{can})}"] \&
\mathbb O_{12}^{*}\times
\underbrace{(\mathfrak S_3)_{S_0}}_{\unsure{(H^+/\mu)_{S_0}|(H^+/\mu_{\mathfrak S})_{S_0}}}
\arrow[r] \& 1
\end{tikzcd}}
it gives, by passage to $(\ )(\widehat{\mathbb Z})^\Delta$, an exact sequence
\diagram{\begin{tikzcd}[ampersand replacement=\&,column sep=small]
(232) \& 1 \arrow[r,dash] \& \widehat M \arrow[r] \& G_{\rho,\sigma}/\mu_G \arrow[r] \&
\widehat{\mathbb Z}^{*}\times\mathfrak S_3 \arrow[r,dash] \& 1,
\end{tikzcd}}
which gives back the operation already considered of $\widehat{\mathbb Z}^{*}\times\mathfrak S_3$
on $\widehat M$, It contains the exact sequence
\diagram{\begin{tikzcd}[ampersand replacement=\&,column sep=small]
(233) \& 1 \arrow[r] \& \widehat M \arrow[r] \&
\begin{gathered}
\inserted{\mathfrak S^{+\wedge}/\mu_{\mathfrak S}}\\
\deleted{\unsure{G_{\rho,\sigma}^{\wedge}/\mu_G|\mathfrak S_{0,3}^{+\wedge}/\mu_{\mathfrak S}}}
\end{gathered}
\arrow[r,very thick] \& \mathfrak S_3 \arrow[r] \& 1
\end{tikzcd}}
which \deleted{\unsure{is}} is a quotient of the \unsure{homotopy} exact sequence

\margin{99}